\chapter{The stack on the robot}
\label{ch:i-robot}

From here on the Mac is on the robot's Wi-Fi (\code{X30*****}, password \code{12345678}), or on
a cable to the Debug Ethernet port at the tail.

\section{SSH keys}

Every host takes user \code{ysc} with password \code{'} (one single quote). Install your key
once, so the commands below never ask again:

\begin{lstlisting}[language=bash]
ssh-copy-id ysc@192.168.1.106
ssh-copy-id ysc@192.168.1.105
\end{lstlisting}

\section{Copy and unpack}

\begin{lstlisting}[language=bash]
cd ~/x30
scp x30_sources.tar.gz ysc@192.168.1.106:~/
scp x30_sources.tar.gz ysc@192.168.1.105:~/
\end{lstlisting}

On each host (\cmd{ssh ysc@192.168.1.106}, then the same on .105):

\begin{lstlisting}[language=bash]
mkdir -p ~/x30_ws && tar -xzf ~/x30_sources.tar.gz -C ~/x30_ws
\end{lstlisting}

The workspace is \file{~/x30_ws} with \file{src/} (x30\_ros and its sources) and
\file{Livox-SDK/}.

\section{Dependencies}
\label{sec:i-apt}

\file{check_deps.sh} reads only what is installed (\cmd{dpkg}, the workspace,
\file{/usr/local/lib}) and prints one \cmd{apt install} line for what is missing:

\begin{lstlisting}[language=bash]
bash ~/x30_ws/src/x30_ros/tools/check_deps.sh                # navigation host
bash ~/x30_ws/src/x30_ros/tools/check_deps.sh --perception   # perception host
\end{lstlisting}

The robot borrows the Mac's internet through SSH for the missing packages. The Mac must then be
on the robot \emph{and} online: a cable to the Debug Ethernet port and Wi-Fi to the internet, or
the robot's Wi-Fi and a phone on USB tethering.

\begin{lstlisting}[language=bash]
ssh -R 1080 ysc@192.168.1.106      # on the Mac: a SOCKS proxy back to it while it is open
# in that session, on the robot:
P='-o Acquire::http::Proxy=socks5h://127.0.0.1:1080 -o Acquire::https::Proxy=socks5h://127.0.0.1:1080'
sudo apt $P update && sudo apt $P install -y <the packages check_deps.sh printed>
\end{lstlisting}

Where SSH forwarding is refused, the robot names the files it needs and the Mac fetches them:

\begin{lstlisting}[language=bash]
sudo apt-get install --print-uris -qq <packages> | cut -d"'" -f2 > ~/urls.txt   # robot
scp ysc@192.168.1.106:~/urls.txt . && mkdir -p debs \
  && (cd debs && xargs -n1 curl -fsSLO < ../urls.txt)                         # Mac, online
scp -r debs ysc@192.168.1.106:~/ && ssh ysc@192.168.1.106 'sudo dpkg -i ~/debs/*.deb'
\end{lstlisting}

\section{Build}

\textbf{Navigation host} (.106): the Livox SDK once, then the whole workspace.

\begin{lstlisting}[language=bash]
cd ~/x30_ws/Livox-SDK/build && cmake .. && make -j"$(nproc)" && sudo make install
cd ~/x30_ws && catkin_make -DCATKIN_ENABLE_TESTING=OFF -DCMAKE_BUILD_TYPE=Release
\end{lstlisting}

\textbf{Perception host} (.105): only the cameras, the recorder and what launches them. Run the
\cmd{catkin_make} line \file{check_deps.sh --perception} prints once everything is there; it
names the packages with \code{-DCATKIN_WHITELIST_PACKAGES}.

\begin{notebox}
A \cmd{pip install --user} that put a setuptools in \file{~/.local} stops catkin
(\code{module 'importlib_metadata' has no attribute 'EntryPoints'}). \file{check_deps.sh}
warns, and its build line then starts with \code{PYTHONNOUSERSITE=1}, which leaves
\file{~/.local} out without touching it.
\end{notebox}

\section{Deep Robotics' autonomy off}
\label{sec:i-dr}

The navigation host also runs Deep Robotics' localization, navigation, mapping, URDF publisher
and its bridge to the robot. Next to the stack they only take CPU and get in its way
(\topic{/initialpose}, \topic{/map}, TF, the UDP port of the motion host's reports). Turn them
off, now and at boot:

\begin{lstlisting}[language=bash]
bash ~/x30_ws/src/x30_ros/tools/dr_services.sh off
bash ~/x30_ws/src/x30_ros/tools/dr_services.sh status   # 'udp 43897 ... free'
\end{lstlisting}

The sensor drivers (\topic{/livox/lidar}, \topic{/imu/data}), the ROS master and the robot
server stay. The motion host and the perception host are never touched. \cmd{dr_services.sh on}
puts everything back as it was.

\section{The motion host's reports}
\label{sec:i-networktoml}

The motion host sends its 200 Hz reports only to the addresses in its \file{network.toml}
(Deep Robotics' Application Manual, 4.3.2). Add the navigation host with port 43897 and restart
the motion program; otherwise the robot takes commands but reports nothing, and the driver's
state says \code{no reports from 192.168.1.103 yet (network.toml on the motion host?)}.

\section{The hosts and their clocks}

Deep Robotics' drivers on the navigation host register under the name \code{ysc-perception}.
On the perception host, make that name resolve:

\begin{lstlisting}[language=bash]
echo '192.168.1.106 ysc-perception' | sudo tee -a /etc/hosts     # on .105
\end{lstlisting}

Without internet the two clocks drift apart; nothing in the stack depends on it, but logs and
reports of the two hosts then disagree. Let the navigation host be the perception host's clock
with \cmd{chrony} (\code{allow 192.168.1.0/24} and \code{local stratum 10} on .106,
\code{server 192.168.1.106 iburst} on .105), or once after each boot on .105:

\begin{lstlisting}[language=bash]
sudo date -s "$(ssh ysc@192.168.1.106 date -Ins)"
\end{lstlisting}
